\textit{Physique I, année 2023 — filière PC}

\section*{III \quad Formation avalancheuse de courant}

L'origine des sylphes est toujours discutée à l'heure actuelle.
Un modèle repose sur la formation d'un courant lié à des ionisations en cascade provoquées par un électron primaire capable d'ioniser par chocs successifs des molécules de l'atmosphère, principalement du diazote $\text{N}_2$.
Les électrons fils créés par l'électron primaire peuvent à leur tour provoquer de nouvelles ionisations d'où l'effet d'avalanche. Cet électron primaire peut provenir du rayonnement cosmique qui éclaire en permanence la terre. Nous allons tout d'abord étudier la possibilité pour un tel électron de provoquer l'ionisation d'une molécule de diazote.

On considère un modèle unidimensionnel d'axe $Ox$ où $N_0$ électrons se déplacent de façon homogène dans le sens $x$ croissant. Ces électrons sont répartis aléatoirement dans une section donnée $\Sigma$. Leur déplacement s'effectue dans un milieu peu dense où la densité volumique de molécules est $n_{\text{mo}}$. On s'intéresse à la distance moyenne $\langle x \rangle$ qu'ils vont parcourir avant de subir un choc avec une molécule. Cette distance porte le nom de libre parcours moyen : $\ell_{\text{pm}} = \langle x \rangle$. On appelle $N(x) < N_0$ le nombre d'électrons qui ont atteint l'abscisse $x$ sans subir de choc depuis $x = 0$. La probabilité qu'un électron subisse un choc entre les abscisses $x$ et $x + \text{d}x$ est le rapport de la section efficace totale associée à l'ensemble des molécules présentes entre $x$ et $x + \text{d}x$ et de la section $\Sigma$. À chaque molécule est associée une section efficace $S_{\text{eff}}$ qui traduit le fait que si l'électron arrive sur cette surface, le choc se produit et s'il passe à côté, il n'y a pas de choc.

\begin{itemize}
    \item[\textbf{18.}] En raisonnant entre $x$ et $x + \text{d}x$, établir l'équation différentielle d'évolution de $N(x)$. Résoudre cette équation et exprimer $N(x)$ en fonction de $N_0$, $n_{\text{mo}}$, $S_{\text{eff}}$ et $x$.
    \item[\textbf{19.}] Montrer que le libre parcours moyen s'écrit $\ell_{\text{pm}} = \langle x \rangle = \frac{1}{n_{\text{mo}} S_{\text{eff}}}$
    \item[\textbf{20.}] Pour la molécule de diazote, on a $S_{\text{eff}} \simeq 10^{-20}\ \text{m}^2$. Justifier cet ordre de grandeur pour la section efficace. Déterminer la valeur numérique du libre parcours moyen $\ell_{\text{pm}}$ pour un électron à l'altitude $z_3 = 60\ \text{km}$ d'altitude où la pression est $P_3 = 20\ \text{Pa}$ et la température $T_3 = 260\ \text{K}$.
\end{itemize}

Au niveau des sylphes, il existe un champ électrique qui a pour principale origine les charges présentes à la surface des nuages d'une part et celles présentes à la limite entre la mésosphère et l'ionosphère vers $100\ \text{km}$ d'altitude. Son intensité est $E \simeq 100\ \text{V}\cdot\text{m}^{-1}$.

\begin{itemize}
    \item[\textbf{21.}] Sachant que l'énergie d'ionisation de la molécule de diazote est $\mathcal{E}_i \simeq 16\ \text{eV}$, un électron initialement au repos et accéléré par le champ électrique local d'intensité $E$ peut-il ioniser cette molécule ? On apportera des arguments numériques.
\end{itemize}

On soupçonne fortement des électrons cosmiques relativistes possédant une énergie cinétique de $1\ \text{MeV}$ d'être à l'origine de la cascade avalancheuse.

\begin{itemize}
    \item[\textbf{22.}] Justifier le fait que ces électrons cosmiques sont nécessairement relativistes.
    \item[\textbf{23.}] Expliquer le phénomène de cascade avalancheuse.
\end{itemize}
